\documentclass[../main.tex]{subfiles}

\begin{document}

\chapter{记号}\label{chapter:appA-记号}

\section{矩阵记号}\label{section:appA-矩阵记号}

\begin{enumerate}[(i)]
\item 我们写 $A=((a_{ij}))$ 表示 $A$ 是一个 $m\times n$ 矩阵，其 $(i,j)$ 元为 $a_{ij}$。（有时，如 \S8.1.4 中，用上标表示行会更方便。）

对角矩阵记为 $\operatorname{diag}(d_1,\cdots,d_n)$。

\item $\mathbb{M}^{m\times n}=$ 实 $m\times n$ 矩阵空间。

$\mathbb{S}^n=$ 实对称 $n\times n$ 矩阵空间。

\item $\operatorname{tr}A=$ 矩阵 $A$ 的迹。

\item $\det A=$ 矩阵 $A$ 的行列式。

\item $\operatorname{cof}A=$ 矩阵 $A$ 的余子式矩阵（见 \S8.1.4）。

\item $A^T=$ 矩阵 $A$ 的转置。

\item 若 $A=((a_{ij}))$ 和 $B=((b_{ij}))$ 是 $m\times n$ 矩阵，则
\begin{align*}
A:B=\sum_{i=1}^m\sum_{j=1}^n a_{ij}b_{ij},
\end{align*}
且
\begin{align*}
|A|=(A:A)^{1/2}=\left(\sum_{i=1}^m\sum_{j=1}^n a_{ij}^2\right)^{1/2}.
\end{align*}

\item 若 $A\in\mathbb{S}^n$，且如下 $x=(x_1,\cdots,x_n)\in\mathbb{R}^n$，相应的二次型为 $x\cdot Ax=\sum_{i,j=1}^n a_{ij}x_ix_j$。

\item 若 $A,B\in\mathbb{S}^n$，我们写
\begin{align*}
A\geqslant B
\end{align*}
表示 $A-B$ 是非负定的。特别地，写
\begin{align*}
A\geqslant\theta I
\end{align*}
若 $x\cdot Ax\geqslant\theta|x|^2$ 对所有 $x\in\mathbb{R}^n$。

\item 有时我们写 $yA$ 表示 $A^Ty$，其中 $A\in\mathbb{M}^{m\times n}$，$y\in\mathbb{R}^m$。
\end{enumerate}

\section{几何记号}\label{section:appA-几何记号}

\begin{enumerate}[(i)]
\item $\mathbb{R}^n=n$ 维实 Euclid 空间，$\mathbb{R}=\mathbb{R}^1$。

$S^{n-1}=\partial B(0,1)=\mathbb{R}^n$ 中的 $(n-1)$ 维单位球面。

\item $e_i=(0,\cdots,0,1,\cdots,0)=$ 第 $i$ 个标准坐标向量。

\item $\mathbb{R}^n$ 中一个典型点为 $x=(x_1,\cdots,x_n)$。根据上下文，我们也将 $x$ 视为行向量或列向量。

\item $\mathbb{R}^n_+=\{x=(x_1,\cdots,x_n)\in\mathbb{R}^n\mid x_n>0\}=$ 开上半空间。

$\mathbb{R}_+=\{x\in\mathbb{R}\mid x>0\}$。

\item $\mathbb{R}^{n+1}$ 中一点常记为 $(x,t)=(x_1,\cdots,x_n,t)$，通常将 $t=x_{n+1}$ 解释为时间。

$x\in\mathbb{R}^n$ 中一点有时写为 $x=(x',x_n)$，其中 $x'=(x_1,\cdots,x_{n-1})\in\mathbb{R}^{n-1}$。

\item $U,V,W$ 通常表示 $\mathbb{R}^n$ 的开子集。我们写
\begin{align*}
V\subset\subset U
\end{align*}
若 $\overline{V}\subset U$ 且 $\overline{V}$ 紧，称 $V$ 紧包含于 $U$ 中。

\item $\partial U=U$ 的边界，$\overline{U}=U\cup\partial U=U$ 的闭包。

\item $U_T=U\times(0,T]$。

\item $\Gamma_T=\overline{U_T}\setminus U_T=U_T$ 的抛物边界。

\item $B(x,r)=\{y\in\mathbb{R}^n\mid|x-y|\leqslant r\}=\mathbb{R}^n$ 中以 $x$ 为心、半径 $r>0$ 的闭球。

\item $B^0(x,r)=$ 以 $x$ 为心、半径 $r>0$ 的开球。

\item $C(x,t;r)=\{y\in\mathbb{R}^n,\,s\in\mathbb{R}\mid|x-y|\leqslant r,\;t-r^2\leqslant s\leqslant t\}=$ 以上中心 $(x,t)$、半径 $r>0$、高 $r^2$ 的闭柱体。

\item $\alpha(n)=\mathbb{R}^n$ 中单位球 $B(0,1)$ 的体积 $=\dfrac{\pi^{n/2}}{\Gamma(n/2+1)}$。

$n\alpha(n)=\mathbb{R}^n$ 中单位球面 $\partial B(0,1)$ 的表面积。

\item 若 $a=(a_1,\cdots,a_n)$ 和 $b=(b_1,\cdots,b_n)$ 属于 $\mathbb{R}^n$，
\begin{align*}
a\cdot b=\sum_{i=1}^n a_ib_i,\quad|a|=\left(\sum_{i=1}^n a_i^2\right)^{1/2}.
\end{align*}

\item $\mathbb{C}^n=n$ 维复空间；$\mathbb{C}=$ 复平面。

若 $z\in\mathbb{C}$，写 $\operatorname{Re}(z)$ 表示 $z$ 的实部，$\operatorname{Im}(z)$ 表示虚部。
\end{enumerate}

\section{函数记号}\label{section:appA-函数记号}

\begin{enumerate}[(i)]
\item 若 $u:U\to\mathbb{R}$，写
\begin{align*}
u(x)=u(x_1,\cdots,x_n) \quad (x\in U).
\end{align*}
称 $u$ 光滑，只要 $u$ 无穷可微。

\item 若 $u,v$ 是两个函数，写
\begin{align*}
u\equiv v
\end{align*}
表示 $u$ 恒等于 $v$；即函数 $u,v$ 对其自变量的所有值一致。我们用记号
\begin{align*}
u\triangleq v
\end{align*}
定义 $u$ 等于 $v$。

函数 $u$ 的支集记为 $\operatorname{spt}u$。

\item $u^+=\max(u,0)$，$u^-=-\min(u,0)$，$u=u^+-u^-$，$|u|=u^++u^-$。

符号函数为
\begin{align*}
\operatorname{sgn}(x)=
\begin{cases}
1 & \text{若 } x>0,\\
0 & \text{若 } x=0,\\
-1 & \text{若 } x<0.
\end{cases}
\end{align*}

\item 若 $\mathbf{u}:U\to\mathbb{R}^m$，写
\begin{align*}
\mathbf{u}(x)=(u_1(x),\cdots,u_m(x)) \quad (x\in U).
\end{align*}
函数 $u_k$ 是 $\mathbf{u}$ 的第 $k$ 个分量，$k=1,\cdots,m$。

\item 若 $\Sigma$ 是 $\mathbb{R}^n$ 中光滑 $(n-1)$ 维曲面，写
\begin{align*}
\int_\Sigma f\,\mathrm{d}S
\end{align*}
表示 $f$ 在 $\Sigma$ 上关于 $(n-1)$ 维曲面测度的积分。若 $C$ 是 $\mathbb{R}^n$ 中曲线，记
\begin{align*}
\int_C f\,\mathrm{d}l
\end{align*}
表示 $f$ 在 $C$ 上关于弧长的积分。

\item 平均：
\begin{align*}
\fint_{B(x,r)}f\,\mathrm{d}y\triangleq\frac{1}{\alpha(n)r^n}\int_{B(x,r)}f\,\mathrm{d}y=f \text{ 在球 } B(x,r) \text{ 上的平均},
\end{align*}
以及
\begin{align*}
\fint_{\partial B(x,r)}f\,\mathrm{d}S\triangleq\frac{1}{n\alpha(n)r^{n-1}}\int_{\partial B(x,r)}f\,\mathrm{d}S=f \text{ 在球面 } \partial B(x,r) \text{ 上的平均}.
\end{align*}
更一般地，
\begin{align*}
\fint_E f\,\mathrm{d}\mu\triangleq\frac{1}{\mu(E)}\int_E f\,\mathrm{d}\mu=f \text{ 在集合 } E \text{ 上的平均},
\end{align*}
只要 $\mu(E)>0$。

\item $\chi_E(x)=\begin{cases}1 & \text{若 } x\in E,\\ 0 & \text{若 } x\notin E;\end{cases}$ $\chi_E$ 是 $E$ 的指示函数。

\item 函数 $u:U\to\mathbb{R}$ 称为 Lipschitz 连续，若
\begin{align*}
|u(x)-u(y)|\leqslant C|x-y|
\end{align*}
对某常数 $C$ 和所有 $x,y\in U$。我们写
\begin{align*}
\operatorname{Lip}[u]\coloneqq\sup_{\substack{x,y\in U\\x\neq y}}\frac{|u(x)-u(y)|}{|x-y|}.
\end{align*}

\item 函数 $f,g$ 的卷积记为 $f*g$。
\end{enumerate}

\subsection{导数记号}\label{subsection:appA-导数记号}

设 $u:U\to\mathbb{R}$，$x\in U$。

\begin{enumerate}[(i)]
\item $\dfrac{\partial u}{\partial x_i}(x)=\lim_{h\to 0}\dfrac{u(x+he_i)-u(x)}{h}$，只要此极限存在。

\item 我们通常写 $u_{x_i}$ 表示 $\dfrac{\partial u}{\partial x_i}$。

\item 类似地，$\dfrac{\partial^2u}{\partial x_i\partial x_j}=u_{x_ix_j}$，$\dfrac{\partial^3u}{\partial x_i\partial x_j\partial x_k}=u_{x_ix_jx_k}$，等等。

\item 多重指标记号：
\begin{enumerate}[(a)]
\item 形如 $\alpha=(\alpha_1,\cdots,\alpha_n)$ 的向量，其中每个分量 $\alpha_i$ 为非负整数，称为阶为 $|\alpha|=\alpha_1+\cdots+\alpha_n$ 的多重指标。

\item 给定多重指标 $\alpha$，定义
\begin{align*}
D^\alpha u(x)\triangleq\frac{\partial^{|\alpha|}u(x)}{\partial x_1^{\alpha_1}\cdots\partial x_n^{\alpha_n}}=\partial_{x_1}^{\alpha_1}\cdots\partial_{x_n}^{\alpha_n}u.
\end{align*}

\item 若 $k$ 为非负整数，$D^ku(x)\coloneqq\{D^\alpha u(x)\mid|\alpha|=k\}$，即所有 $k$ 阶偏导数的集合。对各偏导数赋予某种次序，我们也可将 $D^ku(x)$ 视为 $\mathbb{R}^{n^k}$ 中的点。

\item $|D^ku|=\left(\sum_{|\alpha|=k}|D^\alpha u|^2\right)^{1/2}$。

\item 特殊情形：若 $k=1$，我们将 $Du$ 的元素排列为向量：
\begin{align*}
Du\coloneqq(u_{x_1},\cdots,u_{x_n})=\text{梯度向量}.
\end{align*}
因此 $Du\in\mathbb{R}^n$。我们有时也写
\begin{align*}
u_r\coloneqq\frac{x}{|x|}\cdot Du \quad (x\neq 0)
\end{align*}
表示 $u$ 的径向导数。

若 $k=2$，我们将 $D^2u$ 的元素排列为矩阵：
\begin{align*}
D^2u\coloneqq
\begin{pmatrix}
u_{x_1x_1} & \cdots & u_{x_1x_n}\\
& \ddots & \\
u_{x_nx_1} & \cdots & u_{x_nx_n}
\end{pmatrix}_{n\times n}
=\text{Hessian 矩阵}.
\end{align*}
因此 $D^2u\in\mathbb{S}^n$，即实对称 $n\times n$ 矩阵空间。
\end{enumerate}

\item $u$ 的 Laplace 算子为
\begin{align*}
\Delta u=\sum_{i=1}^n u_{x_ix_i}.
\end{align*}
因此 $\Delta u=\operatorname{div}Du=\operatorname{tr}(D^2u)$。

\item 我们有时在符号 $D,D^2$ 等上附加下标以表示求导变量。例如若 $u=u(x,y)$（$x\in\mathbb{R}^n$，$y\in\mathbb{R}^m$），则 $D_xu=(u_{x_1},\cdots,u_{x_n})$，$D_yu=(u_{y_1},\cdots,u_{y_m})$。
\end{enumerate}

\subsection{函数空间}\label{subsection:appA-函数空间}

\begin{enumerate}[(i)]
\item $C(U)=\{u:U\to\mathbb{R}\mid u\text{ 连续}\}$。

$C(\overline{U})=\{u\in C(U)\mid u\text{ 在 }U\text{ 的有界子集上一致连续}\}$。

$C^k(U)=\{u:U\to\mathbb{R}\mid u\text{ }k\text{ 次连续可微}\}$。

$C^k(\overline{U})=\{u\in C^k(U)\mid D^\alpha u\text{ 在 }U\text{ 的有界子集上一致连续，对所有 }|\alpha|\leqslant k\}$。

因此若 $u\in C^k(\overline{U})$，则对每个多重指标 $\alpha$（$|\alpha|\leqslant k$），$D^\alpha u$ 连续延拓到 $\overline{U}$。

\item $C^\infty(U)=\{u:U\to\mathbb{R}\mid u\text{ 无穷可微}\}=\bigcap_{k=0}^\infty C^k(U)$。

$C^\infty(\overline{U})=\bigcap_{k=0}^\infty C^k(\overline{U})$。

\item $C_c(U),C_c^k(U)$ 等表示 $C(U),C^k(U)$ 等中具有紧支集的函数。

\item $L^p(U)=\{u:U\to\mathbb{R}\mid u\text{ Lebesgue 可测，}\|u\|_{L^p(U)}<\infty\}$，其中
\begin{align*}
\|u\|_{L^p(U)}\triangleq\left(\int_U|u|^p\,\mathrm{d}x\right)^{1/p} \quad (1\leqslant p<\infty).
\end{align*}

$L^\infty(U)=\{u:U\to\mathbb{R}\mid u\text{ Lebesgue 可测，}\|u\|_{L^\infty(U)}<\infty\}$，其中
\begin{align*}
\|u\|_{L^\infty(U)}\triangleq\operatorname*{ess\,sup}_U|u|.
\end{align*}

$L^p_{\mathrm{loc}}(U)=\{u:U\to\mathbb{R}\mid u\in L^p(V)\text{ 对每个 }V\subset\subset U\}$。（另见 \S D.1。）

\item $\|Du\|_{L^p(U)}=\lll Du\rrr_{L^p(U)}$。

$\|D^2u\|_{L^p(U)}=\lll D^2u\rrr_{L^p(U)}$。

\item $W^{k,p}(U),H^k(U)$ 等（$k=0,1,2,\cdots$，$1\leqslant p\leqslant\infty$）表示 Sobolev 空间：见\cref{chapter:ch05-Sobolev空间}。

\item $C^{k,\beta}(\overline{U})$（$k=0,\cdots$，$0<\beta\leqslant 1$）表示 H\"older 空间：见\cref{chapter:ch05-Sobolev空间}。

\item $x$ 和 $t$ 的函数。有时引入在 $x$ 和 $t$ 变量上具有不同光滑性的函数空间是有用的，尽管对此类空间没有标准记号。本书中我们写
\begin{align*}
C^2_1(U_T)=\{u:U_T\to\mathbb{R}\mid u,D_xu,D_x^2u,u_t\in C(U_T)\}.
\end{align*}
特别地，若 $u\in C^2_1(U_T)$，则 $u,D_xu$ 等连续到顶面 $U\times\{t=T\}$。
\end{enumerate}

\section{向量值函数}\label{section:appA-向量值函数}

\begin{enumerate}[(i)]
\item 若 $m>1$ 且 $\mathbf{u}:U\to\mathbb{R}^m$，$\mathbf{u}=(u_1,\cdots,u_m)$，对每个多重指标 $\alpha$ 定义
\begin{align*}
D^\alpha\mathbf{u}=(D^\alpha u_1,\cdots,D^\alpha u_m).
\end{align*}
则
\begin{align*}
D^k\mathbf{u}=\{D^\alpha\mathbf{u}\mid|\alpha|=k\}
\end{align*}
且
\begin{align*}
|D^k\mathbf{u}|\triangleq\left(\sum_{|\alpha|=k}|D^\alpha\mathbf{u}|^2\right)^{1/2},
\end{align*}
如前。

\item 在 $k=1$ 的特殊情形，写
\begin{align*}
D\mathbf{u}\coloneqq
\begin{pmatrix}
u^1_{x_1} & \cdots & u^1_{x_n}\\
\vdots & \ddots & \vdots\\
u^m_{x_1} & \cdots & u^m_{x_n}
\end{pmatrix}_{m\times n}
=\text{梯度矩阵}.
\end{align*}

\item 若 $m=n$，有
\begin{align*}
\operatorname{div}\mathbf{u}\coloneqq\operatorname{tr}(D\mathbf{u})=\sum_{i=1}^n u^i_{x_i}=\mathbf{u}\text{ 的散度}.
\end{align*}

\item 空间 $C(U;\mathbb{R}^m),L^p(U;\mathbb{R}^m)$ 等由那些函数 $\mathbf{u}:U\to\mathbb{R}^m$（$\mathbf{u}=(u_1,\cdots,u_m)$）组成，其中 $u_i\in C(U),L^p(U)$ 等（$i=1,\cdots,m$）。
\end{enumerate}

\begin{remark}\label{remark:appA-上下标说明}
关于上下标的说明。如上所示，我们遵循约定将取值于 $\mathbb{R}^m$（$m>1$）的映射（或取值于 Banach 或 Hilbert 空间的映射）设为粗体。此类映射的分量函数用上标表示。另一方面，$\mathbb{R}^n$ 中典型点 $x$ 不是粗体，分量用下标：$x=(x_1,\cdots,x_n)$。矩阵值映射也设为粗体，其分量函数用上标或根据上下文混用上下标。
\end{remark}

\section{估计记号}\label{section:appA-估计记号}

\begin{definition}\label{definition:appA-常数}
常数。我们用字母 $C$ 表示任何可由已知量显式计算的常数。因此 $C$ 表示的确切值在给定计算中可以逐行变化。最大优点是我们的计算看起来更简洁，因为我们不断将额外因子吸收到项 $C$ 中。
\end{definition}

\begin{definition}\label{definition:appA-大O小O}
\begin{enumerate}[(i)]
\item （大 $O$ 记号。）我们写
\begin{align*}
f=O(g) \quad \text{当 } x\to x_0,
\end{align*}
只要存在常数 $C$ 使得
\begin{align*}
|f(x)|\leqslant C|g(x)|
\end{align*}
对所有充分接近 $x_0$ 的 $x$。

\item （小 $o$ 记号。）我们写
\begin{align*}
f=o(g) \quad \text{当 } x\to x_0,
\end{align*}
只要
\begin{align*}
\lim_{x\to x_0}\frac{|f(x)|}{|g(x)|}=0.
\end{align*}
\end{enumerate}
\end{definition}

\begin{remark}\label{remark:appA-极限说明}
表达式"$O(g)$"（或"$o(g)$"）本身未定义。必须始终伴有极限，例如上面的"当 $x\to x_0$"，尽管此极限常是隐含的。
\end{remark}

\section{关于记号的说明}\label{section:appA-记号说明}

前述记号在 PDE 文献中大体标准，有几个重要例外：

\begin{enumerate}[(i)]
\item 我们使用符号"$Du$"而非"$\nabla u$"表示函数 $u$ 的梯度。原因是"$D^2u$"自然表示 $u$ 的 Hessian 矩阵，而"$\nabla^2u$"会与物理和工程文献中 Laplace 算子的标准记号混淆。多重指标记号用字母 $D$ 看起来也更好。

\item 大多数偏微分方程书籍和论文用"$\Omega$"表示 $\mathbb{R}^n$ 中给定 PDE 成立的开子集。如上所示，我们改为主要使用符号"$U$"表示此类区域。优点有几个。首先，由于典型解记为 $u$，用 $U$ 表示其定义域是合理的，不必换用希腊字母。此外，一旦我们称给定开集为 $U$，字母 $V$ 和 $W$ 就可用于子区域。最后，将 $\Omega$ 保留为概率空间的标准符号很重要。许多重要偏微分方程有概率表示公式（见 \cite{E2}，Freidlin~\cite{Fd} 等）；尽管这些超出本书范围，避免未来记号混淆的可能似乎是明智的。
\end{enumerate}

\end{document}